\section{Verification and reproducibility}
\label{sec:verification}

The proofs above supply the analytic justification.  Computation was
used to check finite algebra, branch and indexing conventions, and
numerical evaluations by independent routes.  It was not used to
replace the contour estimates, continuation arguments, or the
cutoff-to-Laurent comparison.

\subsection{Exact finite algebra}

The accompanying \texttt{code/power\_algebra.py} implements the coefficient
construction and strict-word stuffle recursion with exact rational
symbolic coefficients.  The script \texttt{code/exact\_checks.py}
verifies the finite difference recurrence through power eight, and
compares the harmonic-word expansion against actual finite nested
sums at several distinct indices.  It also checks the full fifth
spectral support, the low-power scalar reductions, the closed residue
polynomial, and the canonical fifth polynomial.  Its output is stored
in \texttt{data/exact\_checks.json}; the reusable word tables are in
\texttt{data/power\_coefficients.json}.

These finite checks have a precise limitation.  For example, verifying
the recurrence at finitely many powers does not prove its assertion
for arbitrary $k$; Proposition~\ref{prop:difference} does.  Likewise,
symbolic reduction of a finite polynomial does not justify exchanging
an infinite sum with an endpoint finite part.  That step is supplied
by Theorems~\ref{thm:ML}, \ref{thm:contour}, and \ref{thm:all-power}.

\subsection{An independent fifth-power calculation}

The script \texttt{code/independent\_fifth.py} imports none of the main
coefficient implementation.  It reconstructs the fifth pole data,
the higher-pole contribution, and the cutoff constant through its own
finite stuffle calculation.  It then evaluates $Q_5$ in two ways.
The first uses local Taylor subtraction of the pole of $\psi(x)^5$,
so no singular numerical quadrature is performed without subtraction.
The second uses the ordinary arithmetic series
\eqref{eq:Q5ordinary}, with an independently generated Bernoulli and
Hurwitz-zeta tail.  With 90 working decimal digits, cutoff $N=96$,
and tail powers through $n^{-36}$, the two routes gave
\begin{equation}
 Q_5\approx
 -8.69030396262453826629558966236565357677029354492388429266.
 \label{eq:q5decimal}
\end{equation}
Their absolute discrepancy was $9.60\times10^{-59}$.
These are floating-point diagnostics, not an interval-certified
enclosure.  The decimal in \eqref{eq:q5decimal} is therefore reported
as a diagnostic approximation; the exact assertion is
Theorem~\ref{thm:q5}.

The ordered-generator checks compare the Gamma formulas with direct
finite nested sums, including nonreal tagged exponents and nonintegral
lower endpoints.  They also check Stieltjes and log-Gamma derivative
consequences, the two-tag staircase, and the polylogarithmic formulas
at negative markers $-1$ and $-2$.  All 23 numerical comparisons in
\texttt{code/verify\_ordered\_generators.py} passed at 70 working
digits; the largest absolute residual was $4.92\times10^{-70}$.
Their defining finite sums and power-series evaluations are separate
from the closed formulas being checked.  These are again
non-interval diagnostics.

\subsection{How to reproduce the package}

The standalone \LaTeX{} file builds with an ordinary \texttt{pdflatex}
installation and the listed standard packages; no external images,
bibliography service, or shell escape is required.  The modular source
fragments are included to facilitate repository integration, while the
single-file source is sufficient for rebuilding the PDF.  Python
requirements are \texttt{sympy} and \texttt{mpmath}; the versions used
for the recorded checks are given in their JSON receipts.  Run
\begin{verbatim}
python3 code/exact_checks.py
python3 code/independent_fifth.py
python3 code/verify_ordered_generators.py
python3 code/build_article.py
\end{verbatim}
from the package directory.  The optional Gaussian relation experiment
has separate instructions and requires a C++17 compiler.

The integration directory contains a claim ledger, the pinned source
revision and hashes, proposed placement notes, and the exact scope of
the corrections.  It deliberately distinguishes proved additions,
recovered antecedents, diagnostics, and unresolved targets.  Nothing
in the package changes the conjecture status of unrelated source
statements.

\section{Further exact research questions}
\label{sec:questions}

The strongest next questions concern finite identities and arithmetic
reduction.  Estimates are useful to justify operations, but are not
the principal intended deliverable.

\subsection{The first mixed spectral coordinate}

Determine whether $\Omega_4$ in \eqref{eq:mixed-def} can be expressed
in a specified algebra generated by ordinary zeta derivatives,
ordinary Stieltjes constants, and explicitly named Gamma or Barnes
periods.  Equations~\eqref{eq:OmegaT}--\eqref{eq:Tpolylog} give
three concrete starting points: nonsymmetric derivatives of double
zeta functions, a Tornheim derivative, and a dilogarithm integral.
An answer should identify which independent-looking coordinates are
actually removed.  Merely renaming the dilogarithm integral does not
constitute an arithmetic reduction; a failed integer-relation search
does not prove nonreducibility.

One tractable first task is to study identities for
$\partial_1Z(2,2)$ and $(\partial_1-\partial_2)Z(3,1)$ together,
instead of trying to evaluate $D_{1,3}$ in isolation.  Differentiated
stuffle determines symmetric combinations, whereas the second
expression is explicitly antisymmetric in the spectral directions.
A new identity must address that remaining direction.

\subsection{Scalar polynomials and the growth of spectral support}

The normalized fifth polynomial in
Corollary~\ref{powers:cor:canonical-fifth} cancels the displayed
lower height-one coefficients, but retains $\Omega_4$.  For general
$k$, determine the image of scalar polynomials
$P(\psi)$ under the finite-word spectral map of
Theorem~\ref{thm:algorithm}.  One can first ask a wholly finite
algebraic question: in a stated presentation modulo specified stuffle
and differential relations, which coordinates can be eliminated by
lower powers?

A useful all-order result would give a recurrence for normalized
monic polynomials and a proved description of their remaining word
support.  The stronger assertion that only the single highest
height-one Laurent coefficient remains is not established by the
degree-four cancellation and is not established here at degree five.
An arithmetic simplification of the mixed coordinate could change
the conclusion; formal coordinates must not be silently promoted to
independent numbers.

\subsection{Explicit sixth and seventh moments}

The included coefficient tables through power eight make the next
cases finite, but a table is not yet a well-organized theorem.
Carry out $Q_6$ and $Q_7$ with a reduced, explicitly chosen spectral
coordinate family.  The all-power formula already supplies the
boundary constants $C_6=-\pi^6/7$ and $C_7=0$, and the complete
Stieltjes subtraction is explicit.  The remaining task is to organize
the ordinary multiple-zeta values in $J_k$ and the word derivatives
in $D_k$ without losing diagonal terms.

Compare the resulting support with the fifth mixed coordinate.
Does a finite preferred family emerge at each fixed weight after
all differentiated stuffle identities are imposed?  The deliverable
should include exact coordinate changes and reproducible rational
certificates, not only a numerical basis.

\subsection{A common Hurwitz shift with correct endpoint terms}

For $a>0$, seek an analogue for
$\int_0^1\psi(a+x)^k\,\dd x$ in a common-shift multiple-Hurwitz
notation.  A derivation must retain both ordinary endpoints
$a$ and $a+1$.  Then study the singular limit $a\downarrow0$
using the original coordinate and compare it with $Q_k$.
Replacing $\gamma_j$ by $\gamma_j(a)$ in the unshifted formula
without deriving the new pole locations and endpoint contributions
does not answer this question.

The relative functions in Section~\ref{sec:ordered} suggest a useful
language for endpoint differences.  Their spectral indices are
entire, but that alone does not determine the finite part of a
shifted digamma product.  A precise theorem should specify the
shift in every denominator and prove the connection.

\subsection{Differential polynomials in digamma and polygamma}

Extend the method from $\psi^k$ to an arbitrary fixed differential
polynomial
$P(\psi,\psi',\ldots,\psi^{(r)})$ with constant coefficients.
Its local pole data are finite and its large-circle growth is
controlled, suggesting a generalized normalized expansion and a
finite-part reduction.  Derivatives may change the angular
constant substantially, so it should be computed rather than
assumed to be the same $C_k$.

A first exact target is the family
$\FP\int_0^1\psi(x)^a\psi'(x)^b\,\dd x$ with integers
$a,b\ge0$.  Integration by parts gives relations between these
moments, and the pole-coefficient method supplies their boundary
terms.  Classifying the independent moments after these identities
would connect the primitive formulas to a broader polygamma calculus.

\subsection{Several tags at independent endpoints}

Corollary~\ref{ord:staircase} is proved only for finite intervals.
For independent positive endpoints $a,b$, construct the full
upper-triangular matrix generator with several tagged exponents,
prove normal convergence in all depth markers, and determine
whether the integer staircase still collapses to a shifted ordered
relative sum.  A direct matrix and asymptotic subtraction argument
would be preferable to interpolation from integer intervals, which
does not by itself prove the extension.

If such a theorem holds, differentiating at several spectral poles
should produce exact combinations of multivariate harmonic Laurent
data.  The useful target is an explicitly evaluated mixed derivative
combination, with its continuation independent of the order of
taking removable limits.

\subsection{Boundary markers and removable Gamma singularities}

The one-marker theorem provides an ordinary depth sum at $x=-m$
when $a>m$.  Study the boundary $a=m$ and the exceptional points
$a=m-\ell$ through the actual generating function, separating
ordinary convergence, meromorphic continuation, and finite part.
At finite intervals the generating function is a polynomial and
some apparent Gamma poles are removable; for independent endpoints
this cancellation needs a proof.  Determine the finite limiting
identities and their log-Gamma or Stieltjes corrections in each
admissible case.

\subsection{Polylogarithmic cancellation at roots of unity}

The identities of Proposition~\ref{ord:polylog} hold normally on
$|z|<1$.  Analyze their boundary continuation at $z=1,-1,i,-i$,
with explicit conditions on the tagged exponent.  At $z=1$ the
prefactors and divergent depth sums interact, so a boundary identity
must be derived with a specified limiting operation.  At other roots
of unity, seek exact colored multiple-zeta formulas and their
spectral derivatives.  A rigorous boundary theorem may supply
relations useful to the Gaussian part of the manuscript.

\subsection{Gaussian relations beyond the searched family}

The computation in Appendix~\ref{app:s6} shows only that one
explicitly generated finite family does not close the current
$S_6$ target modulo one prime.  A productive next step is to derive
new exact functional equations or construct the full relevant
Cayley ideal at a higher filtration, then seek a rational
certificate for the target.  Increasing a numerical precision or
repeating the same finite-field search does not supply a proof.
The revised $S_8$ target requires its own exact treatment and must
not inherit a proved label from a fitted expression.

\subsection{Formal verification of the analytic normalization}

The finite coefficient recurrence and the strict-word conversion
are suitable for small proof-assistant components.  The most
valuable analytic components would formalize the full-circle
digamma estimate, normal convergence of the subtracted
Mittag--Leffler series, and the cutoff-to-Laurent lemma.
Together they would make the all-power theorem reusable without
re-proving an endpoint regularization separately at each degree.
An exact algebra certificate alone should not be advertised as
formal verification of the analytic theorem.

\appendix
\section{A bounded investigation of the current Gaussian \texorpdfstring{$S_6$}{S6} target}
\label{app:s6}

For clarity, the searched target is recorded independently of the
proved digamma results.  In this appendix use the colored convention
\[
 \Li_{a,b}(z_1,z_2)=
       \sum_{n>m\ge1}\frac{z_1^nz_2^m}{n^am^b},\qquad
 g_{a,b}=\Im\Li_{a,b}(i,1),
\]
and $\beta(s)=\sum_{n\ge0}(-1)^n/(2n+1)^s$.
The current short residual in the pinned manuscript is
\begin{align}
 R_6={}&S_6+\frac{12334}{527}g_{6,1}
          +\frac{384}{31}g_{5,2}+\frac{2136}{527}g_{4,3}
 \notag\\
 &-\frac{3179}{5713920}\pi^7
     +\frac{801}{2108}\beta(4)\zeta(3)+2\beta(6)\log2,
 \label{eq:S6residual}\\
 S_6={}&\sum_{n\ge0}\frac{(-1)^nH_n}{(2n+1)^6}.
 \notag
\end{align}
The conjecture is $R_6=0$; no such equality is proved here.

We extended the explicitly generated weight-seven relation family
to maximum total input depth five.  After imaginary projection,
primitive integer normalization, and duplicate removal, the family
had the following size:
\begin{center}
\begin{tabular}{lr}
\toprule
Relation family & Rows\\
\midrule
Double-shuffle relations & $14\,908$\\
Distribution relations and specified product lifts & $4\,938$\\
Lifts of the proved $S_2/S_4$ relations & $128$\\
Balanced Cayley relations and specified lifts & $2\,322$\\
\midrule
Total & $22\,296$\\
\bottomrule
\end{tabular}
\end{center}
There were $8\,786$ occupied imaginary word coordinates.  A completed
sparse elimination over $\mathbb F_{2147483647}$ found rank
$8\,705$ and left a nonzero $16$-coordinate representative of the
target.  The family was exhausted in approximately $328$ seconds
on the machine used, before the stipulated $600$-second cap.
Timing is not a mathematical invariant; the field, row family,
rank, and nonzero reduced target are the reproducible diagnostic.

The imaginary evaluation of the exact integer target is $128588R_6$.
A separate rational-arithmetic script verifies this normalization,
including the conversion $\beta(7)=61\pi^7/184320$.
The family generator, field eliminator, complete-run receipts,
normalization check, and reproduction note are included in
\texttt{exploratory/S6\_depth5\_bounded\_search.zip}.
The raw matrix is regenerated by the included code.

This establishes nonmembership in the \emph{specified finite-field
row span}.  It does not establish rational-span nonmembership:
an integer matrix can lose rank modulo the selected prime.
It does not establish nonvanishing of $R_6$, and it does not prove
the conjecture.  Moreover, the finite balanced Cayley family here
is not claimed to generate the full Cayley shuffle ideal at depth
five.  This distinction prevents comparison with stronger exact
obstruction statements in the incoming material from being
overstated.  The retained log is a reproducible computation record,
not a separately replayed compact rational certificate.

\subsection{A useful local attribution}

The canonical $S_2$ section would benefit from a citation immediately
beside its classical formula
\begin{equation}
 S_2=2\Im\Li_3\!\left(\frac{1+i}{2}\right)
       -G\log2-\frac{\pi(\log2)^2}{16}-\frac{\pi^3}{64},
 \qquad G=\beta(2).
 \label{eq:S2attribution}
\end{equation}
Xu and Wang, Section~3.4, Example~3.13 of the inspected arXiv
version~3 \cite{XuWang}, give the identical evaluation after
$R_{1,\bar2}=4S_2$ is substituted.  The broader manuscript already
cites that work; the proposed change is a local attribution
improvement, not a correction to the value or its proof.

\begin{thebibliography}{99}
\addcontentsline{toc}{section}{References}
\raggedright

\bibitem{Repo}
V. Reshetnikov and contributors,
\emph{ProveIt}, source revision
\texttt{b7667ae40a2957eed9a53d8559d53a3ee20dd1c1},
inspected 11 October 2026.
Manuscript and incoming reports:
\url{https://github.com/VladimirReshetnikov/ProveIt/tree/b7667ae40a2957eed9a53d8559d53a3ee20dd1c1/Analysis/Polylogarithms/docs/manuscript}
and
\url{https://github.com/VladimirReshetnikov/ProveIt/tree/b7667ae40a2957eed9a53d8559d53a3ee20dd1c1/docs/incoming}.

\bibitem{Quartic}
ProveIt research continuation,
\emph{Quartic Digamma Bridges, Harmonic Tail Moments, and Complex
Resolvent Powers}, 11 October 2026.
Incoming archive
\texttt{ProveIt\_Quartic\_Tails\_Resolvents\_2026-10-11.zip}
at the revision in \cite{Repo}, especially the quartic bridge
and further-research Questions~1--4.

\bibitem{Barnes}
ProveIt research continuation,
\emph{Barnes Resolvents and Higher Harmonic Identities},
11 October 2026.
Incoming archive
\texttt{ProveIt\_Barnes\_Resolvents\_and\_Harmonic\_Identities\_2026-10-11.zip}
at the revision in \cite{Repo}, especially the quartic harmonic
section and its normalization.

\bibitem{Independent}
ProveIt research continuation,
\emph{Independent Orders and Exact Reductions in
Polylogarithm--Zeta Calculus}, 11 October 2026.
Incoming archive
\texttt{ProveIt\_Independent\_Orders\_2026-10-11.zip}
at the revision in \cite{Repo}, especially the relative
harmonic interpolation and exact Gaussian relation framework.

\bibitem{DLMF}
NIST Digital Library of Mathematical Functions.
Sections~5.5 (functional relations), 5.7 (series expansions),
5.11 (asymptotics and Gamma ratios), 25.11 (Hurwitz zeta),
25.12 (polylogarithms), and 25.16 (multiple zeta values).
\url{https://dlmf.nist.gov/}.

\bibitem{IKZ}
K. Ihara, M. Kaneko, and D. Zagier,
Derivation and double shuffle relations for multiple zeta values,
\emph{Compositio Mathematica} \textbf{142} (2006), 307--338.
\url{https://doi.org/10.1112/S0010437X0500182X}.

\bibitem{INY}
K. Ihara, Y. Nakamura, and S. Yamamoto,
Interpolant of truncated multiple zeta functions.
arXiv:2407.20509 (2024); published in
\emph{The Ramanujan Journal} \textbf{67} (2025), Article~1.
\url{https://arxiv.org/abs/2407.20509}.
\url{https://doi.org/10.1007/s11139-025-01045-2}.
Proposition~2.1 and Corollary~4.2 are cited using the arXiv text.

\bibitem{Young}
P. T. Young,
Series of height one multiple zeta functions,
\emph{Integers} \textbf{24} (2024), A43.
\url{https://math.colgate.edu/~integers/y43/y43.pdf}.
\url{https://doi.org/10.5281/zenodo.11221606}.

\bibitem{BB2014}
D. H. Bailey and J. M. Borwein,
\emph{Computation and theory of Mordell--Tornheim--Witten sums II}.
Author-hosted version dated 5 May 2014, inspected at
\url{https://www.davidhbailey.com/dhbpapers/mtw2.pdf}.
The correction in Section~\ref{powers:sec:audit} is specifically to
equation~(110), printed page~24, in this version.

\bibitem{Mezo2014}
I. Mez\H{o}, Nonlinear Euler sums,
\emph{Pacific Journal of Mathematics} \textbf{272} (2014), no.~1,
201--226.
\url{https://doi.org/10.2140/pjm.2014.272.201}.
\url{https://msp.org/pjm/2014/272-1/pjm-v272-n1-p09-s.pdf}.
The wording discussed in Section~\ref{powers:sec:audit} is on
printed page~205.

\bibitem{XuWang}
C. Xu and W. Wang,
\emph{Dirichlet type extensions of Euler sums}.
arXiv:2009.11704v3, 12 April 2022,
Section~3.4, Example~3.13.
\url{https://arxiv.org/html/2009.11704v3}.

\end{thebibliography}
\end{document}
